\documentclass[11pt]{article}
\usepackage[margin=1.15in]{geometry}
\usepackage{amsmath,amssymb,amsthm,mathtools}
\usepackage{xcolor}
\usepackage[colorlinks=true,linkcolor=blue!60!black,citecolor=blue!60!black,urlcolor=blue!60!black]{hyperref}

\newtheorem{theorem}{Theorem}[section]
\newtheorem{proposition}[theorem]{Proposition}
\newtheorem{lemma}[theorem]{Lemma}
\newtheorem{corollary}[theorem]{Corollary}
\theoremstyle{definition}
\newtheorem{definition}[theorem]{Definition}
\theoremstyle{remark}
\newtheorem{remark}[theorem]{Remark}

\newcommand{\C}{\mathbb{C}}
\newcommand{\R}{\mathbb{R}}
\newcommand{\Z}{\mathbb{Z}}
\newcommand{\w}{\omega}
\newcommand{\Sp}{\mathrm{Sp}}
\newcommand{\Stab}{\mathrm{Stab}}
\newcommand{\tr}{\operatorname{tr}}
\newcommand{\diag}{\operatorname{diag}}
\newcommand{\dg}{\dagger}
\newcommand{\arccosh}{\operatorname{arccosh}}

\title{The Pais--Uhlenbeck Metric as a Minimum-Distortion Principle,\\
and the Representation Problem for the Fourth-Order Field}
\author{}
\date{}

\begin{document}
\maketitle

\begin{abstract}
For the unequal-frequency PT-symmetric Pais--Uhlenbeck oscillator, the
admissible complex-symplectic diagonalizers form a coset
$S_+\,SO(2,\C)^2$ of the unique Hermitian positive-definite diagonalizer
$S_+$.  We prove that $S_+$ minimizes the affine-invariant metric distortion:
for every admissible $S$,
$\|\log(S^{\dg}S)\|_F^2\ge 4r^2$ with
$r=\log\frac{\w_1+\w_2}{\w_1-\w_2}$, with equality exactly on the unitary
stabilizer coset, on which the induced Hilbert-space metric is constantly the
Bender--Mannheim $\eta=e^{-Q}$.  The proof is elementary and exact: at
worst-case stabilizer alignment the log-eigenvalues of $S^{\dg}S$ are
$\arccosh(\cosh r\cosh b)\pm a$ in the Gram rapidities $(a,b)$ of the
stabilizer factor.  We then quantize the free fourth-order field modewise and
compute the Bender--Mannheim ground state exactly: it is a normalizable real
Gaussian for every mode, but its fidelity with the two-field Fock vacuum
tends to the universal constant $\sqrt3/2$ at large momentum, and its
relative particle content to $1/3$ per mode pair.  Consequently the physical
(PT) and naive Fock representations of the free fourth-order field are
disjoint in every spatial dimension, although each mode is unitarily
equivalent: PT unitarization is exact modewise and changes the
field-theoretic representation.  All identities are machine-verified.
\end{abstract}

\section{Introduction}\label{sec:intro}

This paper continues the symplectic analysis of the PT-symmetric
Pais--Uhlenbeck (PU) oscillator begun in \cite{companion}, whose results and
conventions we use freely: $V=(x,y,p,q)^{\top}$, $[V_a,V_b]=iJ_{ab}$,
$H_{\mathrm{PT}}=\tfrac12V^{\top}GV$ with the PU matrix $G$, positive normal
form $G_0$, rapidity
\begin{equation}
r=\log\frac{\w_1+\w_2}{\w_1-\w_2},\qquad \w_1>\w_2>0,
\end{equation}
and, after the canonical normalization $d_xd_y=\gamma\w_1\w_2$, the unique
Hermitian positive-definite diagonalizer
\begin{equation}
S_+=\cosh\tfrac r2\,I+\sinh\tfrac r2\,B\in\Sp(4,\C),
\qquad
S_+^{\top}G'S_+=G_0',\qquad S_+^{\top}JS_+=J,
\end{equation}
with $B$ a fixed Hermitian involution.  The full solution family of the
diagonalization problem is the coset $S_+\,\Stab(J,G_0')$ with
$\Stab(J,G_0')\cong SO(2,\C)\times SO(2,\C)$, and the metaplectic logarithm
of $S_+$ is the Bender--Mannheim generator $Q=\alpha pq+\beta xy$
\cite{BM2008PRL,BM2008PRD,companion}.

Two questions are addressed here.

\emph{First}, uniqueness of the positive diagonalizer does not by itself
distinguish $S_+$ variationally: one may ask whether $S_+$ \emph{minimizes} a
natural functional over all admissible diagonalizers.  We prove that it
minimizes the affine-invariant distortion
$\mathcal F(S)=d(I,S^{\dg}S)^2=\|\log(S^{\dg}S)\|_F^2$
(Theorem~\ref{thm:mindist}), so that the Bender--Mannheim metric is selected
by a minimum-distortion principle: it is the least-distorting complex
canonical transformation converting the PU Hamiltonian to its positive
normal form.

\emph{Second}, for the free fourth-order field
$(\Box+m_1^2)(\Box+m_2^2)\phi=0$ each spatial momentum $k$ contributes a PU
pair with $\w_i(k)=\sqrt{k^2+m_i^2}$, and the modewise construction
applies with rapidity
\begin{equation}\label{eq:rk}
r(k)=\log\frac{\big(\w_1(k)+\w_2(k)\big)^2}{m_1^2-m_2^2}
=2\log|k|+O(1)\qquad(m_1>m_2).
\end{equation}
One might expect from \eqref{eq:rk} that the modewise transformation carries
unboundedly growing squeezing and fails Fock implementability
``dramatically''.  The exact computation shows something sharper and more
structured (Theorems~\ref{thm:groundstate}--\ref{thm:disjoint}): the PT
ground state exists as a normalizable Gaussian for \emph{every} mode, its
mismatch with the free two-field vacuum saturates at \emph{universal
constants} --- fidelity $\sqrt3/2$ and relative occupation $1/3$ per mode
pair at large $k$ --- and therefore the physical and naive representations
are disjoint in every spatial dimension $d\ge1$, with vacuum overlap decaying
like $e^{-c V\Lambda^d}$ under a cutoff.  The divergence of $r(k)$ itself is
a statement about the observable-space (phase-space) metric, which we
separate cleanly (Corollary~\ref{cor:scale}): its eigenvalues
$e^{\pm r(k)}\asymp(1+k^2)^{\pm1}$ make the phase-space metric an
$H^1\oplus H^{-1}$-type Sobolev pair, unboundedly inequivalent to the
standard norm, while the state-level mismatch remains $O(1)$ per mode.

\paragraph{Verification.}
As in \cite{companion}, every finite-dimensional identity in this paper ---
the invariant formulas, the closed-form log-eigenvalues, the PT ground-state
Gaussian, the eigenvalue equation, the fidelity and occupation formulas, and
the limits $\sqrt3/2$ and $1/3$ --- is machine-verified symbolically, with
independent numerical checks (global optimization for
Theorem~\ref{thm:mindist}; high-precision quadrature for
Theorem~\ref{thm:groundstate}).  The scripts accompany the verification
repository of \cite{companion}.

\section{The minimum-distortion theorem}\label{sec:mindist}

Work in the canonically normalized coordinates of \cite{companion}; all
adjoints refer to them.  The stabilizer is parametrized by two complex
angles,
\begin{equation}
C(\theta_1,\theta_2)=\sum_{j=1,2}\big(\cos\theta_j\,P_j+\sin\theta_j\,X_j\big),
\qquad X_j^2=-P_j ,
\end{equation}
with $P_j$ the projections onto the modes of $G_0'$ and $X_j$ the per-mode
complex structures; every admissible diagonalizer is $S=S_+C$.

\begin{definition}
For an admissible $S$ let
$\mathcal F(S)=\|\log(S^{\dg}S)\|_F^2$, the squared affine-invariant distance
of the observable-space metric $S^{\dg}S$ from the identity
\cite{Bhatia2007}.  For a stabilizer element $C=C_1\oplus C_2$ (blocks on the
two modes) define its \emph{Gram rapidities} $\tau_j\ge0$ by
\begin{equation}
W_j=C_jC_j^{\dg},\qquad \det W_j=1,\qquad
\cosh\tau_j=\tfrac12\tr W_j ,
\end{equation}
and set $a=\tfrac12(\tau_1+\tau_2)$, $b=\tfrac12(\tau_1-\tau_2)$.
\end{definition}

Note $\tau_1=\tau_2=0$ iff $W_1=W_2=I$ iff $C$ is unitary.  Since
$P=S^{\dg}S=C^{\dg}S_+^2C$ is Hermitian positive \emph{and} symplectic, its
eigenvalues form reciprocal pairs $e^{\pm x_1},e^{\pm x_2}$ with
$x_i\ge0$, and $\mathcal F(S)=2(x_1^2+x_2^2)$.

\begin{lemma}[Invariants]\label{lem:invariants}
Write $W=CC^{\dg}=W_1\oplus W_2$, $c=\cosh r$, $s=\sinh r$.  Decomposing
$W_j=\cosh\tau_j\,I+\sinh\tau_j\,(n_j\!\cdot\!\sigma)$ with real unit vectors
$n_j$, and letting $\chi\in[-1,1]$ denote the alignment
$n_1\!\cdot\!\tilde n_2$ (where $\tilde n_2$ is $n_2$ with its $\sigma_2$
component reflected),
\begin{align}
\cosh x_1+\cosh x_2 &= c\,(\cosh\tau_1+\cosh\tau_2), \label{eq:T}\\
\cosh x_1\cosh x_2 &= \tfrac12\Big[(2c^2-s^2)\cosh\tau_1\cosh\tau_2
+s^2\big(1+\sinh\tau_1\sinh\tau_2\,\chi\big)\Big]. \label{eq:Pi}
\end{align}
\end{lemma}

\begin{proof}
$\tr P=\tr(e^{rB}W)=c\,\tr W+s\,\tr(BW)$, and $\tr(BW)=0$ because $B$ is
off-diagonal and $W$ diagonal with respect to the mode splitting; this gives
\eqref{eq:T}.  For \eqref{eq:Pi}, expand
$\tr P^2=c^2\tr W^2+s^2\tr\big((BW)^2\big)$ (the cross terms again vanish
blockwise); the off-diagonal blocks of $B$ equal $-\sigma_2$, and the $2\times2$
identity $\sigma_2M\sigma_2=(\det M)\,M^{-\top}$ reduces
$\tr\big((BW)^2\big)=2\tr(\overline{W_2}^{\,-1}W_1)
=4\big[\cosh\tau_1\cosh\tau_2-\sinh\tau_1\sinh\tau_2\,\chi\big]$.
Converting $(\tr P,\tr P^2)$ to the symmetric functions of
$\cosh x_i$ yields \eqref{eq:Pi}.  (Both matrix identities are verified
against random stabilizer elements in the repository.)
\end{proof}

\begin{lemma}[Worst alignment; closed form]\label{lem:closedform}
At $\chi=-1$ the system \eqref{eq:T}--\eqref{eq:Pi} is solved exactly by
\begin{equation}\label{eq:xformula}
x_{1,2}=\varphi(b)\pm a,
\qquad
\varphi(b)=\arccosh\big(\cosh r\,\cosh b\big)\ \ge\ r .
\end{equation}
\end{lemma}

\begin{proof}
With $\alpha=\cosh a$, $\beta=\cosh b$ one finds from
\eqref{eq:T}--\eqref{eq:Pi} at $\chi=-1$, using
$\cosh\tau_1\cosh\tau_2\pm\sinh\tau_1\sinh\tau_2=\cosh(\tau_1\pm\tau_2)$,
\[
\cosh x_1+\cosh x_2=2c\,\alpha\beta,
\qquad
\cosh x_1\cosh x_2=\alpha^2+c^2\beta^2-1 ,
\]
whence
$\cosh x_{1,2}=c\alpha\beta\pm\sqrt{(\alpha^2-1)(c^2\beta^2-1)}
=\cosh(\varphi\pm a)$ with $\cosh\varphi=c\beta$.  Both hyperbolic identities
are verified symbolically.  Finally
$\cosh\varphi(b)-\cosh r=\cosh r(\cosh b-1)\ge0$ gives $\varphi(b)\ge r$ with
equality iff $b=0$.
\end{proof}

\begin{lemma}[Alignment monotonicity]\label{lem:monotone}
At fixed $\cosh x_1+\cosh x_2$, the quantity $x_1^2+x_2^2$ is nondecreasing
in $\cosh x_1\cosh x_2$; and \eqref{eq:Pi} is nondecreasing in $\chi$.
\end{lemma}

\begin{proof}
Parametrize $\cosh x_{1,2}=T/2\pm w$, $w\ge0$; the product decreases in $w$,
and
$\frac{d}{dw}(x_1^2+x_2^2)=2\big(\frac{x_1}{\sinh x_1}-\frac{x_2}{\sinh x_2}\big)
\le0$
because $x/\sinh x$ is strictly decreasing and $x_1\ge x_2$.  The second
statement is the sign of the coefficient
$s^2\sinh\tau_1\sinh\tau_2/2\ge0$ in \eqref{eq:Pi}.
\end{proof}

\begin{theorem}[Minimum distortion]\label{thm:mindist}
For every admissible diagonalizer $S=S_+C$,
\begin{equation}
\mathcal F(S)\ \ge\
4\Big[a^2+\arccosh^2\big(\cosh r\cosh b\big)\Big]\ \ge\ 4r^2=\mathcal F(S_+),
\end{equation}
with overall equality if and only if $C$ is unitary.  All minimizers share
the polar factor $S_+$, hence the same Hilbert-space metric
$\eta=e^{-Q}$: the Bender--Mannheim metric is the unique metric of any
minimum-distortion diagonalizer,
\begin{equation}
S_+\in\operatorname*{arg\,min}_{S^{\top}JS=J,\ S^{\top}G'S=G_0'}
d\big(I,S^{\dg}S\big),
\qquad
Q=-2\log\big(\text{polar factor of any minimizer}\big)\big|_{\text{metaplectic}} .
\end{equation}
\end{theorem}

\begin{proof}
By Lemma~\ref{lem:monotone}, lowering $\chi$ to $-1$ at fixed Gram rapidities
does not increase $x_1^2+x_2^2$; by Lemma~\ref{lem:closedform} the value there
is $\mathcal F=2\big[(\varphi+a)^2+(\varphi-a)^2\big]=4(\varphi^2+a^2)$;
and $\varphi(b)\ge r$.  Equality in the chain
forces $a=0$ and $b=0$, i.e.\ $\tau_1=\tau_2=0$, i.e.\ $C$ unitary; the
converse is immediate since then $S^{\dg}S$ is unitarily conjugate to
$S_+^2$.  If $C$ is unitary, $S_+C=S_+\cdot C$ is already a polar
decomposition, so the polar factor is $S_+$ and the metaplectic implementer
factorizes as $\rho'=U_C\,e^{-Q/2}$ with $U_C$ unitary, giving
$\eta'=\rho'^{\dg}\rho'=e^{-Q}$ for every minimizer.
\end{proof}

\begin{remark}
(i) The individual log-eigenvalues do \emph{not} satisfy $x_i\ge r$: numerical
scans show $\min_i x_i<r$ on open regions of the stabilizer; only the joint
quantity is bounded below, which is why the two-invariant argument is needed.
(ii) The equality set is the unitary part of the stabilizer, which in the
canonical coordinates is generically the finite group $\Z_2\times\Z_2$
\cite{companion}; abstractly it is contained in the maximal compact
$U(1)^2$ of $SO(2,\C)^2$.
(iii) The theorem admits a geometric recognition.  $\mathcal F(S)=2\|\mu(S)\|^2$
with $\mu$ the Cartan projection of $\Sp(4,\C)$, and $\mu(S_+)=(r,r)$ lies on
the $C_2$ chamber wall.  Although the stabilizer itself is not
$\theta$-compatible, it embeds in the compatible subgroup
$SL(2,\C)\times SL(2,\C)$ of mode-preserving symplectics, whose orbit through
the basepoint is a totally geodesic $\mathbb H^3\times\mathbb H^3$ inside the
symmetric space $X=\Sp(4,\C)/\Sp(2)$; the Gram data $(\tau_j,n_j)$ are polar
coordinates on the two hyperbolic factors.  Since $B$ is trace-orthogonal to
all mode-preserving Hermitian directions, the $S_+$-geodesic leaves this
orbit orthogonally, and Theorem~\ref{thm:mindist} becomes the nearest-point
(orthogonal projection) theorem for convex totally geodesic subsets of an
NPC space --- including the equality set, by uniqueness of the projection
foot.  The closed form \eqref{eq:xformula} is the corresponding Pythagorean
law: hyperbolic ($\arccosh(\cosh r\cosh b)$, the hypotenuse law of a totally
geodesic $\mathbb H^2$) in $b$, flat in $a$.  Numerical tests confirm the
general principle: for any Hermitian Hamiltonian $\xi$ orthogonal to the
mode-preserving directions, $\|\mu(e^{\xi}C)\|\ge\|\mu(e^{\xi})\|$ on the
stabilizer coset, while generic non-orthogonal $\xi$ violate it.
\end{remark}

\section{The free fourth-order field}\label{sec:field}

\subsection{Modewise data and the phase-space Hilbert scale}

For $(\Box+m_1^2)(\Box+m_2^2)\phi=0$, each spatial momentum $k$ carries a PU
pair with $\w_i(k)=\sqrt{k^2+m_i^2}$ and $m_1>m_2>0$.  Since
$(\w_1-\w_2)(\w_1+\w_2)=m_1^2-m_2^2$ identically, the rapidity obeys
\eqref{eq:rk} \emph{exactly}, interpolating between
$r(0)=\log\frac{m_1+m_2}{m_1-m_2}$ and $r(k)=2\log|k|+O(1)$.

\begin{corollary}[Phase-space Hilbert scale]\label{cor:scale}
In the canonically normalized modewise coordinates, the observable-space
metric $M_{\mathrm{obs}}(k)=S_+(k)^{\dg}S_+(k)$ has eigenvalues
$e^{\pm r(k)}$ with
\begin{equation}
e^{r(k)}=\frac{\big(\w_1(k)+\w_2(k)\big)^2}{m_1^2-m_2^2}
\asymp 1+k^2 ,
\qquad
e^{-r(k)}\asymp(1+k^2)^{-1} .
\end{equation}
On the spectral subspaces $\Pi_\pm(k)$ of $B$ the induced one-particle
phase-space norm is therefore equivalent to an $H^1\oplus H^{-1}$ Sobolev
pair (up to the constant $m_1^2-m_2^2$ and bounded factors): it is not
uniformly equivalent to the standard norm, and no scalar renormalization
makes it so.
\end{corollary}

This is the precise content of the observation that the modewise metric
``blows up in the UV''.  It concerns the classical (phase-space) metric.
The quantum state-level comparison is finer, and is computed next.

\subsection{The PT ground state, exactly}

Fix a mode and work in the normalized coordinates $(x,y)$ with
$\w_1>\w_2>0$.  The positive Hamiltonian $\widetilde H'$ has ground state
$\phi_0\propto e^{-\frac12(x^2/\w_2+\w_1y^2)}$.  The PT Hamiltonian in these
coordinates is the differential operator
\begin{equation}
H_{\mathrm{PT}}'=\tfrac12\Big[\tfrac{\w_1^2+\w_2^2}{\w_1\w_2}\,x^2
+\w_1\w_2\,y^2-\w_1\w_2\,\partial_x^2\Big]-x\,\partial_y .
\end{equation}

\begin{theorem}[PT ground state]\label{thm:groundstate}
The unique Gaussian solution of the modewise annihilation equations
$\big(\ell_j^{\top}S_+^{-1}V\big)\psi_0=0$ is
\begin{equation}
\psi_0(x,y)\ \propto\
\exp\!\Big[-\tfrac12\Big(\tfrac{\w_1+\w_2}{\w_1\w_2}\,x^2
+2xy+(\w_1+\w_2)\,y^2\Big)\Big],
\end{equation}
a \emph{real} Gaussian with
$\det(\mathrm{quadratic\ form})=\frac{\w_1^2+\w_1\w_2+\w_2^2}{\w_1\w_2}>0$:
normalizable for all $\w_1>\w_2>0$, and
$H_{\mathrm{PT}}'\psi_0=\tfrac{\w_1+\w_2}{2}\,\psi_0$.  Its fidelity with the
two-oscillator vacuum and its occupation relative to it are exactly
\begin{align}
f(\w_1,\w_2)&=\big|\langle\hat\phi_0|\hat\psi_0\rangle\big|^2
=\frac{4\w_1\sqrt{\w_1^2+\w_1\w_2+\w_2^2}}{4\w_1^2+3\w_1\w_2+\w_2^2},\\[2pt]
\langle N\rangle(\w_1,\w_2)
&=\langle\hat\psi_0|(N_1+N_2)|\hat\psi_0\rangle
=\frac{\w_2(\w_1^2+\w_2^2)}{2\w_1(\w_1^2+\w_1\w_2+\w_2^2)},
\qquad \langle N_1\rangle=\langle N_2\rangle .
\end{align}
In particular, along the field dispersion $\w_i=\w_i(k)$,
\begin{equation}\label{eq:universal}
\lim_{|k|\to\infty}f=\frac{\sqrt3}{2}\approx0.866,
\qquad
\lim_{|k|\to\infty}\langle N\rangle=\frac13 ,
\end{equation}
universal constants independent of $m_1,m_2,\gamma$.
\end{theorem}

\begin{proof}
Solving the two first-order equations for a Gaussian ansatz gives the stated
quadratic form (one solution); the eigenvalue equation is a direct
substitution.  The fidelity is the standard two-Gaussian determinant formula;
the occupations follow from covariance algebra,
$\langle N_j\rangle=\tfrac12 c_j^{\top}\Sigma c_j$ with
$\Sigma$ the covariance of $|\psi_0|^2$ and $c_j$ the coefficient vectors of
$a_j\psi_0$.  All steps are verified symbolically, and the fidelity is
confirmed by 25-digit quadrature.  The limits follow with
$\w_1/\w_2\to1$.
\end{proof}

\begin{remark}
Two features deserve emphasis.  (i) $\psi_0$ is a smooth, normalizable
Gaussian even as $\w_1\to\w_2$: the exceptional-point pathology
(divergence of $S_+$, Jordan structure) lives in the similarity
transformation and the excited spectrum, not in the ground state itself.
(ii) A naive Bogoliubov estimate would assign the modewise map a squeezing
$\sim\sinh^2\!\big(r(k)/2\big)\sim k^2$ and hence a violently divergent
particle content per mode.  The exact result \eqref{eq:universal} is
$O(1)$: the diverging rapidity $r(k)$ is carried almost entirely by
directions that do not create quanta.  The naive estimate fails because
$\rho=e^{-Q/2}$ is not $*$-preserving on the original Fock structure ---
the transformed ``annihilators'' $\rho a_j\rho^{-1}$ do not satisfy the
unitary Bogoliubov normalization ($\alpha\alpha^{\dg}-\beta\beta^{\dg}
=\cosh r\ne1$) --- so Shale--Stinespring-type criteria cannot be applied to
$\rho$ directly.  Comparing the two \emph{states}, both of which live in the
same $L^2$, avoids this trap entirely.
\end{remark}

\subsection{Disjointness of the representations}

\begin{theorem}[Representation change]\label{thm:disjoint}
Let $m_1>m_2>0$, $d\ge1$, and consider the free fourth-order field in a box
of volume $V$ with momentum cutoff $\Lambda$, with
$\Psi_{\mathrm{PT}}=\bigotimes_k\hat\psi_0(k)$ the modewise PT ground state
and $\Omega=\bigotimes_k\hat\phi_0(k)$ the free two-field Fock vacuum.  Then:
\begin{enumerate}
\item the vacuum overlap satisfies
$\big|\langle\Omega|\Psi_{\mathrm{PT}}\rangle\big|^2
=\prod_k f(k)\le e^{-c\,V\Lambda^d}$ for some $c>0$ depending only on
$m_1/m_2$ and $d$ (since $f(k)\to\sqrt3/2<1$);
\item the relative particle density diverges in the UV:
$\frac1V\langle N_{\mathrm{tot}}\rangle
=\int^{\Lambda}\!\frac{d^dk}{(2\pi)^d}\,\langle N\rangle(k)
\sim\frac13\int^{\Lambda}\!\frac{d^dk}{(2\pi)^d}=\infty$ as
$\Lambda\to\infty$, in \emph{every} dimension $d\ge1$;
\item in the thermodynamic/continuum limit the incomplete-tensor-product
classes of $\Psi_{\mathrm{PT}}$ and $\Omega$ differ
($\sum_k(1-|\langle\hat\psi_0|\hat\phi_0\rangle|)=\infty$), so the GNS
representations of the CCR algebra built on the two states are disjoint: no
density matrix in the Fock representation of the free pair
$(m_1,m_2)$ reproduces the PT vacuum, and vice versa.
\end{enumerate}
Each mode, by contrast, is unitarily equivalent: the modewise Dyson maps
$\rho_k=e^{-Q_k/2}$ define unitaries between the mode Hilbert spaces
equipped with the physical ($\eta_k$) and standard inner products.
\end{theorem}

\begin{proof}
(1) and (2) are immediate from Theorem~\ref{thm:groundstate}: $f(k)$ and
$\langle N\rangle(k)$ are continuous in $k$ with the universal limits
\eqref{eq:universal}, so beyond some $k_*(m_1/m_2)$ one has, e.g.,
$f(k)\le0.9$ and $\langle N\rangle(k)\ge0.3$.  (3) is von Neumann's
equivalence criterion for incomplete tensor products \cite{vN1939} applied
to the two product vectors, combined with purity: inequivalent pure product
states generate disjoint representations.  The modewise statement is the
finite-dimensional similarity of \cite{BM2008PRD,companion}.
\end{proof}

\begin{remark}[What this does and does not say]\label{rem:meaning}
The theorem does not say the PT quantization of the fourth-order field is
inconsistent.  The physical Hilbert space, defined as the completion of the
modewise-unitarized theory, is naturally identified with the Fock space of
the free pair $(m_1,m_2)$; on it the dynamics is unitary, exactly as argued
by Bender and Mannheim modewise.  What fails is the \emph{global} equivalence
between this physical representation and the naive covariant Fock
representation in which the fourth-order propagator and perturbation theory
are usually written: the two are different, mutually singular
representations of the same field algebra, with an $O(1)$ discrepancy per UV
mode --- not a discrepancy that decays with $k$, and not one confined to the
exceptional locus $m_1=m_2$.  Any interacting construction therefore has to
choose a representation first; formal manipulations that mix them (e.g.,
inserting $\eta=e^{-Q}$ as if it were a bounded operator on the covariant
Fock space) are not supported: by Corollary~\ref{cor:scale} the metric is
unbounded with unbounded inverse, defining a Hilbert scale rather than an
equivalent norm, and by Theorem~\ref{thm:disjoint} its formal
infinite-mode product annihilates the covariant vacuum's representation
class entirely.  This sharpens, for the free theory, the sense in which
PT-based unitarity claims for quadratic gravity
\cite{QQG2024} require an explicit domain-and-completion construction.
\end{remark}

\section{Discussion}\label{sec:discussion}

\paragraph{The variational principle.}
Theorem~\ref{thm:mindist} upgrades the uniqueness theorem of
\cite{companion}: the Bender--Mannheim metric is not merely the unique
positive point of the diagonalizer coset but its unique distortion minimizer,
with the exact excess $4[a^2+\arccosh^2(\cosh r\cosh b)]-4r^2$ in the Gram
rapidities of the stabilizer.  The result is a concrete, fully explicit
instance of the Kempf--Ness/Mostow picture --- minimal-norm points on
complexified orbits are unitary-locus points --- but the proof is elementary
and yields the minimizing value in closed form, which the abstract theory
does not.

\paragraph{The two geometries, again.}
The field-theoretic analysis separates two statements that are easily
conflated.  The \emph{phase-space} metric degenerates in the UV exactly as
the rapidity formula \eqref{eq:rk} dictates ($e^{\pm r(k)}\asymp k^{\pm2}$, a
Sobolev pair $H^1\oplus H^{-1}$); this is a statement about observables.
The \emph{state-space} mismatch saturates at universal constants
($\sqrt3/2$, $1/3$); this is a statement about representations, and it is
what controls implementability.  Both are exact consequences of the same
modewise data.  The constants are universal because the UV limit
$\w_1/\w_2\to1$ erases all mass and coupling dependence --- remarkably, the
same limit in which the finite-dimensional theory degenerates (the
exceptional point) is the limit in which the ground-state data become
featureless.

\paragraph{Outlook.}
Three natural continuations: (i) extend the minimum-distortion theorem to
general semisimple quadratic PT Hamiltonians, where the stabilizer is a
product of $SO(2,\C)$ factors and the same Gram-rapidity reduction should
apply mode-pair by mode-pair; (ii) determine whether the disjointness
survives interactions in low dimensions, where the free-theory obstruction
is only the universal $1/3$ per mode (a wave-function renormalization of the
modewise map, $\rho_k\to\rho_k Z_k$ with $Z_k$ unitary, cannot change class
(3), but a $k$-dependent redefinition of the mode split might); and (iii)
make contact with the Krein-space formulation of indefinite-metric QFT,
where the covariant representation lives, to phrase the representation
choice intrinsically.  We leave these to future work.

\begin{thebibliography}{9}

\bibitem{companion}
\emph{Canonical positive symplectic diagonalization of the Pais--Uhlenbeck
oscillator}, companion paper and verification repository, 2026.

\bibitem{BM2008PRL}
C.~M.~Bender and P.~D.~Mannheim,
Phys.\ Rev.\ Lett.\ \textbf{100} (2008) 110402.

\bibitem{BM2008PRD}
C.~M.~Bender and P.~D.~Mannheim,
Phys.\ Rev.\ D \textbf{78} (2008) 025022.

\bibitem{Bhatia2007}
R.~Bhatia, \emph{Positive Definite Matrices}, Princeton University Press,
2007.

\bibitem{vN1939}
J.~von Neumann,
\emph{On infinite direct products},
Compos.\ Math.\ \textbf{6} (1939) 1.

\bibitem{QQG2024}
J.~Kubo and J.~Kuntz,
\emph{Unitarity through PT symmetry in quantum quadratic gravity},
arXiv:2410.08278.

\end{thebibliography}

\end{document}
